\author{OTS}
\documentclass[12pt]{../../../inc/mybib}

\setincpath{../../../inc/}

\usepackage{bible_style}
\usepackage{textmakro}

\graphicspath{{../../../assets/images/}}


% Absatzformatierung
\setlength{\parindent}{0pt}
\setlength{\parskip}{0.9em}
\renewcommand{\arraystretch}{1.8}
\begin{document}

\begin{longtable}{l p{14cm}}
    %\hline
    % \rowcolor{gray!20}

    \textbf{Vers} & \textbf{Einteilung Nehemia}                                                                                      \\

    \hline
    \endfirsthead

    \hline
    % \rowcolor{gray!20}
    \textbf{Vers} & \textbf{Einteilung Nehemia}                                                                                      \\
    \hline
    \endhead
    \verslink{Nehemia}{1}{1} - \verslink{Nehemia}{1}{4}
                  & Nehemia bekommt Informationen von Jerusalem
    \\
    \hline
    \verslink{Nehemia}{1}{5} - \verslink{Nehemia}{1}{11}
                  & Gebet Nehemias zum \herr N
    \\
    \hline
    \verslink{Nehemia}{2}{1} - \verslink{Nehemia}{2}{6}
                  & Nehemia spricht vor dem König und bittet darum, nach Jerusalem zu gehen
    \\
    \hline
    \verslink{Nehemia}{2}{7} - \verslink{Nehemia}{2}{8}
                  & Vorbereitung für die Reise nach Jerusalem
    \\
    \hline
    \verslink{Nehemia}{2}{9} - \verslink{Nehemia}{2}{10}
                  & Nehemias Reise nach Jerusalem
    \\
    \hline
    \verslink{Nehemia}{2}{11} - \verslink{Nehemia}{2}{16}
                  & Nehemia untersucht die Stadtmauer
    \\
    \hline
    \verslink{Nehemia}{2}{17} - \verslink{Nehemia}{2}{20}
                  & Nehemia gibt das Resultat seiner Untersuchung weiter
    \\
    \hline
    \verslink{Nehemia}{3}{1} - \verslink{Nehemia}{3}{32}
                  & Beginn des Wiederaufbaus und Einteilung der Arbeiter
    \\
    \hline
    \verslink{Nehemia}{3}{33} - \verslink{Nehemia}{4}{17}
                  & Anfechtungen wärend dem Bau von den Widersachern
    \\
    \hline
    \verslink{Nehemia}{5}{1} - \verslink{Nehemia}{5}{13}
                  & Aufstand des Volkes wärend dem Bau
    \\
    \hline
    \verslink{Nehemia}{5}{14} - \verslink{Nehemia}{5}{19}
                  & Nehemia wird Stadthalter von Jerusalem
    \\
    \hline
    \verslink{Nehemia}{6}{1} - \verslink{Nehemia}{6}{14}
                  & Nehemia wird von den Feinden verleumdet
    \\
    \hline
    \verslink{Nehemia}{6}{15} - \verslink{Nehemia}{7}{4}
                  & Der Bau ist vollendet
    \\
    \hline
    \verslink{Nehemia}{7}{5} - \verslink{Nehemia}{7}{72}
                  & Nehemia liest das Geschlechtsregister der zuerst Heraufgezogenen vor
    \\
    \hline
    \verslink{Nehemia}{8}{1} - \verslink{Nehemia}{8}{12}
                  & Esra liest das Gesetzbuch vor
    \\
    \hline
    \verslink{Nehemia}{8}{13} - \verslink{Nehemia}{8}{18}
                  & Feier des Laubhüttenfestes
    \\
    \hline
    \verslink{Nehemia}{9}{1} - \verslink{Nehemia}{9}{4}
                  & Das Volk tut Busse
    \\
    \hline
    \verslink{Nehemia}{9}{5} - \verslink{Nehemia}{9}{37}
                  & Gebet und Lobpreis zu Gott von den Leviten
    \\
    \hline
    \verslink{Nehemia}{10}{1} - \verslink{Nehemia}{10}{40}
                  & Das Volk geht eine neue Verpflichtung ein
    \\
    \hline
    \verslink{Nehemia}{11}{1} - \verslink{Nehemia}{12}{26}
                  & Besiedelung Jerusalem und Umgebung
    \\
    \hline
    \verslink{Nehemia}{12}{27} - \verslink{Nehemia}{12}{42}
                  & Einweihung der neuen Mauer
    \\
    \hline
    \verslink{Nehemia}{13}{1} - \verslink{Nehemia}{13}{6}
                  & Korruption in Abwesenheit Nehemias
    \\
    \hline
    \verslink{Nehemia}{13}{7} - \verslink{Nehemia}{13}{30}
                  & Rückkehr Nehemias und der Versuch die Ordnung wieder herzustellen.
    \\
    \hline
\end{longtable}

\end{document}